\documentclass[12pt]{article}
\usepackage[utf8]{inputenc}
\usepackage[T1]{fontenc}
\usepackage[margin=1in]{geometry}
\usepackage{setspace}
\usepackage{xcolor}
\usepackage[colorlinks=true, urlcolor=blue!60!black, linkcolor=blue!60!black]{hyperref}
\setlength{\parindent}{0em}
\setlength{\parskip}{0.6em}
\renewcommand{\baselinestretch}{1.0}
\begin{document}
\begin{center}
{\LARGE \textit{Cruzan v.\ Director, Missouri Department of Health}, 497 U.S.\ 261 (1990)\par}
\vspace{0.5em}
{\large Case Brief --- Two-Page Edition\par}
\end{center}

\section{Memory Jogger}
A competent person may refuse lifesaving treatment, but a State may require ``clear and convincing evidence'' of an incompetent patient's wishes before life support is withdrawn (\href{https://supreme.justia.com/cases/federal/us/497/261/}{Justia}).

\section{Detailed Case Facts}
After a 1983 car accident, 25-year-old Nancy Cruzan was left in a persistent vegetative state, kept alive by a gastrostomy feeding tube in a Missouri state hospital. Her parents asked that artificial nutrition and hydration be ended; the hospital refused without court approval.

\section{Procedural History}
The trial court authorized withdrawal based on Nancy's statements to a housemate. The Missouri Supreme Court reversed, 4--3, finding the statements ``unreliable'' and requiring clear and convincing evidence. The U.S.\ Supreme Court affirmed.

\section{Judicial Votes}
Affirmed 5--4. Rehnquist wrote for White, O'Connor, Scalia, and Kennedy; O'Connor and Scalia concurred; Brennan (with Marshall and Blackmun) and Stevens dissented.

\section{Holding}
A competent person has a due process liberty interest in refusing unwanted treatment, assumed to include artificial nutrition and hydration. A State may require clear and convincing evidence of an incompetent patient's wishes and need not defer to family ``substituted judgment.''

\section{Analysis of Opinions}
\textbf{Majority:} the State's unqualified interest in life, the risk of family abuse, and the non-adversarial nature of these hearings justify heightened proof, since an erroneous withdrawal is irrevocable. \textbf{O'Connor:} the right to refuse tube feeding is protected, and honoring a designated proxy ``may well be constitutionally required.'' \textbf{Scalia:} ``the federal courts have no business in this field.'' \textbf{Brennan:} the right is fundamental and error runs both ways. \textbf{Stevens:} ``Choices about death touch the core of liberty.''

\section{Examples: Future Applications}
\textbf{Same-Side:} a State may bar ventilator withdrawal absent clear proof of the patient's wishes. \textbf{Opposite-Side:} a hospital may not force-feed a competent patient who refuses. \textbf{Fence-Sitter:} a formally designated proxy with no specific instructions, a question footnote 12 left open.

\section{Critique}
Critics say tube feeding is an ongoing intervention, not a neutral ``status quo,'' and that the rule favors those who draft living wills. Originalists, following Scalia, object that the Court assumed an unenumerated right at all.

\section{Key Quotations}
\begin{quote}
``The principle that a competent person has a constitutionally protected liberty interest in refusing unwanted medical treatment may be inferred from our prior decisions.''
\end{quote}
\noindent --- Chief Justice Rehnquist, for the Court, 497 U.S.\ at 278.

\end{document}
